\documentclass{article}
\usepackage{amsmath, amssymb, amsthm}

\setlength{\parindent}{0pt}

\newtheorem*{claim}{Claim}

\begin{document}

\begin{claim}
    Let $R$ be an equivalence relation on $A$. The equivalence classes of $R$ form a partition of $A$.
\end{claim}

\begin{proof}
    We show that every equivalence class of $R$ is non-empty, every element of $A$ belongs to an equivalence class,
    and any two distinct equivalence classes are disjoint.

    \medskip

    \textbf{Non-empty and coverage.} Let $a \in A$. Since $R$ is reflexive, $aRa$, so $a \in [a]$. Hence every 
    equivalence class is nonempty and every element of $A$ belongs to at least one equivalence class.

    \medskip

    \textbf{Disjoint.} Suppose $[a] \cap [b] \neq \emptyset$. Then there exists $c \in A$ such that $c \in [a]$
    and $c \in [b]$. So $cRa$ and $cRb$. Since $R$ is symmetric, $cRa$ implies $aRc$. Since $R$ is transitive, 
    $aRc$ and $cRb$ implies $aRb$.

    \medskip
    
    If $x \in [a]$, then $aRx$. Since $R$ is symmetric, $aRx$ implies $xRa$.
    Since $R$ is transitive, $xRa$ and $aRb$ implies $xRb$, so $x \in [b]$. Hence $[a] \subseteq [b]$. Similarly, 
    $[b] \subseteq [a]$. Thus $[a] = [b]$.

    \medskip

    Therefore the equivalence classes of $R$ form a partition of $A$.
\end{proof}

\end{document}
